
\section{Implementation of Multi-Suction Cup Gripper Optimization}

This chapter presents the optimization framework developed for a multi-suction cup gripper
consisting of three suction cups arranged in a triangular configuration. The framework supports both
evolutionary multi-objective optimization and Bayesian multi-objective optimization. The
evolutionary algorithms are implemented using the \texttt{pymoo} library~\cite{pymoo}, whereas
Bayesian optimization is implemented using the \texttt{Ax} platform and its underlying
\texttt{BoTorch} modeling framework~\cite{olson2025ax,botorch}.

For both optimization approaches, the evaluation of candidate gripper layouts is delegated to an
external service through an API. This separation provides a modular architecture in which the
optimization strategy and the domain-specific grasp evaluation remain decoupled. Consequently, the
same decision-variable representation, grasp-generation procedure, and objective functions can be
used by both evolutionary and Bayesian optimization methods, allowing their performance to be
compared under a common evaluation pipeline.

\subsection{Reduced Parameterization of Multi-Suction Cup Gripper}

The gripper consists of three suction cups whose positions must ultimately be represented in
Cartesian space in order to perform grasp generation, collision checking, and objective computation.
However, the Cartesian suction cup coordinates are not used directly as optimization variables.
Instead, the layout is parameterized in a reduced form to constrain the search space and encode
geometric structure.

All three suction cups are constrained to lie on a common circle in the gripper plane. Their
configuration is therefore defined by one shared radius $r$ and three angular parameters
$\alpha_1,\alpha_2,\alpha_3$:
\begin{equation}
    \mathbf{q} = [r,\alpha_1,\alpha_2,\alpha_3]
\end{equation}
\nomenclature{$\mathbf{q}$}{Reduced decision variable vector defining the suction cup layout (see also $\mathbf{x}$ for generic decision variables)}
\nomenclature{$r$}{Common radial distance of the suction cups from the gripper center}
\nomenclature{$\alpha_i$}{Angular position of the $i$-th suction cup on the circle}

The angular parameters determine the positions of the suction cups on the circle, while the common
radius determines their distance from the gripper center. For a fixed gripper height $z_0$, the
Cartesian coordinates of the suction cups are obtained through the polar-to-Cartesian transformation
\begin{equation}
    \mathbf{c}_i =
    \begin{bmatrix}
        r\cos(\alpha_i) \\
        r\sin(\alpha_i) \\
        z_0
    \end{bmatrix},
    \qquad i=1,2,3.
\end{equation}
\nomenclature{$\mathbf{c}_i$}{Cartesian position of the $i$-th suction cup in the gripper frame}
\nomenclature{$z_0$}{Fixed height coordinate of the suction cups}

The resulting Cartesian gripper configuration is written as
\begin{equation}
    \mathbf{p} = [\mathbf{c}_1,\mathbf{c}_2,\mathbf{c}_3].
\end{equation}
\nomenclature{$\mathbf{p}$}{Cartesian suction cup configuration derived from the reduced parameterization (see also $\mathbf{x}$ for generic decision variables)}

Thus, the optimization is carried out in the reduced 4-dimensional parameter space $\mathbf{q}$,
while the grasp evaluation operates on the derived Cartesian coordinates $\mathbf{p}$.

This formulation reduces the dimensionality of the search space from six planar degrees of freedom,
\begin{equation}
    [(x_1,y_1),(x_2,y_2),(x_3,y_3)] \in \mathbb{R}^6,
\end{equation}
to four variables,
\begin{equation}
    [r,\alpha_1,\alpha_2,\alpha_3].
\end{equation}
The reduction improves search efficiency and excludes many geometrically redundant or undesirable
layouts while preserving sufficient flexibility for diverse triangular suction cup arrangements.

\paragraph{Assumption}
A triangular configuration of three suction cups is assumed to yield a stable
grasp~\cite{madan2022realrobotchallenge2021}. This is motivated by the fact that three spatially
distributed contact points form a statically stable support polygon, enabling effective load
distribution and resistance to external disturbances.

\subsection{Decision Variables and Bounds for Multi-Suction Cup Gripper}

The optimization is performed over the reduced decision vector
\begin{equation}
    \mathbf{q} = [r,\alpha_1,\alpha_2,\alpha_3].
\end{equation}

All variables are treated as continuous. The angular variables are bounded by
\begin{equation}
    \alpha_i \in [0,2\pi), \qquad i=1,2,3,
\end{equation}
and the common radius is bounded by
\begin{equation}
    r \in [r_{\min},r_{\max}].
\end{equation}
\nomenclature{$r_{\min}$}{Lower bound on the common suction cup radius (see also $\mathcal{X}$ for decision space bounds)}
\nomenclature{$r_{\max}$}{Upper bound on the common suction cup radius (see also $\mathcal{X}$ for decision space bounds)}

The radial bounds are determined separately for each workpiece. The workpiece is oriented such that its longest dimension lies in the horizontal
plane, the horizontal bounding box is centred at \((x,y)=(0,0)\), and the base of the workpiece is placed at \(z=0\). This establishes a consistent
coordinate system for defining the gripper-design bounds. Let \(X_{\min}\), \(X_{\max}\), \(Y_{\min}\), and \(Y_{\max}\) denote the coordinates of the
bounding box of a workpiece in the \(x\)- and \(y\)-directions. The maximum radial distance is defined by the distance from the origin to the nearest
bounding-box boundary:

\begin{equation}
    r_{\max}
    =
    \min
    \left\{
        |X_{\min}|,\,
        |X_{\max}|,\,
        |Y_{\min}|,\,
        |Y_{\max}|
    \right\}.
\end{equation}


The lower radial bound is derived from the suction-cup diameter
\(d_{\mathrm{cup}}\) and is defined as

\begin{equation}
    r_{\min}
    =
    \frac{\sqrt{3}}{3}
    d_{\mathrm{cup}}.
\end{equation}

This lower bound corresponds to the circumradius of an equilateral triangle with side length \(d_{\mathrm{cup}}\).

All optimization methods generate candidate values within these bounds. In the evolutionary
algorithms, the initial population is obtained by random sampling and subsequently modified through
the operators of the selected algorithm. In Bayesian optimization, Ax first evaluates the center of
the search space and then uses Sobol sampling to generate an initial space-filling design. After
sufficient initialization data have been collected, Ax transitions to model-based Bayesian
optimization.

\subsection{Optimization Framework of Multi-Suction Cup Gripper}

The multi-suction cup gripper problem is solved using either an evolutionary multi-objective
optimization algorithm or Bayesian multi-objective optimization. Both approaches operate on the same
reduced decision vector, variable bounds, objective functions, and external evaluation service.

The common optimization configuration contains:
\begin{itemize}
    \item Number of objectives $n_{\text{obj}}$
    \item Decision-variable definitions and bounds
    \item Optimization direction (minimization or maximization)
    \item Choice of optimization method
    \item Evaluation endpoint URL
\end{itemize}

For evolutionary optimization, the additional configuration parameters include:
\begin{itemize}
    \item Population size $N$
    \item Maximum number of generations $G_{\max}$
    \item Choice of evolutionary algorithm
    \item Algorithm-specific parameters, such as crossover and mutation settings
\end{itemize}

For Bayesian optimization, the additional configuration parameters include:
\begin{itemize}
    \item Number of optimization batches $B_{\max}$
    \item Maximum number of candidates per batch $b$
    \item Bayesian generation strategy
\end{itemize}

In the current Bayesian implementation, the configuration value named \texttt{generations}
represents the number of candidate-generation batches rather than evolutionary generations. The
maximum requested number of Bayesian evaluations is therefore
\begin{equation}
    E_{\max}^{\mathrm{BO}} = B_{\max}b.
\end{equation}
The actual number may be smaller if Ax returns fewer than $b$ candidates in a particular batch.
\nomenclature{$B_{\max}$}{Maximum number of Bayesian optimization batches (see also $B_{\max}$ in two-finger optimization)}
\nomenclature{$b$}{Maximum number of Bayesian candidates requested per batch (see also $b$ in two-finger optimization)}
\nomenclature{$E_{\max}^{\mathrm{BO}}$}{Maximum requested number of Bayesian objective evaluations (see also $E_{\max}^{\mathrm{BO}}$ in two-finger optimization)}

\subsubsection{Supported Optimization Methods}

The framework supports the following evolutionary and Bayesian multi-objective optimization methods:

\begin{table}[H]
\centering
\caption{Supported multi-objective optimization methods.}
\label{tab:algorithms_sc}
\begin{tabularx}{\textwidth}{@{}l l >{\raggedright\arraybackslash}X@{}}
\toprule
\textbf{Family} & \textbf{Identifier} & \textbf{Method} \\
\midrule
Evolutionary & NSGA2 & Non-dominated Sorting Genetic Algorithm II \\
Evolutionary & NSGA3 & Non-dominated Sorting Genetic Algorithm III \\
Evolutionary & MOEAD & Multi-Objective Evolutionary Algorithm based on Decomposition \\
Evolutionary & CMOPSO & Competitive Multi-Objective Particle Swarm Optimization \\
Evolutionary & SMSEMOA & S-Metric Selection Evolutionary Multi-Objective Algorithm \\
Bayesian & AX-MBM & Ax modular BoTorch model for multi-objective Bayesian optimization \\
\bottomrule
\end{tabularx}
\end{table}

% For evolutionary algorithms requiring reference directions, such as NSGA-III and MOEA/D, uniform
% reference directions are generated using the method of Das and Dennis~\cite{das_dennis}:
% \begin{equation}
%     \Lambda =
%     \texttt{get\_reference\_directions}
%     (\text{``das-dennis''},\; n_{\text{obj}},\; n_{\text{partitions}}).
% \end{equation}
% \nomenclature{$\Lambda$}{Set of reference directions for decomposition-based algorithms (see also $\Lambda$ in two-finger optimization)}
% \nomenclature{$n_{\text{obj}}$}{Number of objectives (see also $n_{\text{obj}}$ in two-finger optimization)}

% The Bayesian implementation uses the default Ax generation strategy identified at runtime as
% \texttt{Center+Sobol+MBM:fast}. This strategy consists of three stages:
% \begin{enumerate}
%     \item Evaluation of the center of the bounded search space
%     \item Sobol quasi-random sampling for initialization
%     \item Model-based optimization using the Ax modular BoTorch model (MBM)
% \end{enumerate}

% The Center and Sobol stages do not employ a surrogate model or acquisition function. After
% transitioning to the MBM stage, Ax fits probabilistic surrogate models to the completed observations
% and selects new candidates using an acquisition function appropriate for the multi-objective
% problem. For continuous multi-objective problems, this generally involves Gaussian-process
% surrogates and a hypervolume-improvement acquisition function. The exact model and
% acquisition-function classes are selected by Ax according to the installed version and experiment
% configuration and should therefore be recorded from the runtime configuration to ensure
% reproducibility.

\subsection{Initialization of Multi-Suction Cup Gripper Optimization}

The initialization procedure depends on the selected optimization approach.

\subsubsection{Evolutionary Initialization}

For evolutionary optimization, an initial population of $N$ candidate solutions is generated by
sampling the reduced decision variables within their bounds:
\begin{equation}
    \mathcal{U}^{(0)}
    =
    \{\mathbf{q}_1,\mathbf{q}_2,\dots,\mathbf{q}_N\}.
\end{equation}
\nomenclature{$\mathcal{U}^{(0)}$}{Initial evolutionary population of candidate decision vectors}
\nomenclature{$N$}{Population size (see also $N$ in two-finger optimization)}

Each candidate decision vector $\mathbf{q}_k$ is transformed into its corresponding Cartesian
gripper configuration $\mathbf{p}_k$ before objective evaluation.

\subsubsection{Bayesian Initialization}

Bayesian optimization does not maintain an evolutionary population. Instead, it stores a dataset of
completed trials:
\begin{equation}
    \mathcal{D}_t =
    \left\{
        \left(\mathbf{q}_k,\mathbf{g}(\mathbf{q}_k)\right)
    \right\}_{k=1}^{t},
\end{equation}
where $\mathbf{q}_k$ is an evaluated gripper parameterization and $\mathbf{g}(\mathbf{q}_k)$ is its
observed objective vector. \nomenclature{$\mathcal{D}_t$}{Dataset of completed Bayesian optimization trials after $t$ evaluations (see also $\mathcal{D}_t$ in two-finger optimization)}

The Ax generation strategy first evaluates the center of the search space and then generates an
initial set of candidates using Sobol quasi-random sampling. Sobol sampling provides broad coverage
of the bounded decision space before a surrogate model is fitted. Once the initialization
requirements of the generation strategy have been satisfied, Ax transitions to model-based candidate
generation.

\subsection{Objective Function Evaluation of Multi-Suction Cup Gripper via API}

The evaluation of objective values is decoupled from the optimization loop and delegated to an
external evaluation service. For each candidate solution $\mathbf{q}$, the corresponding Cartesian
suction cup coordinates are first computed by transforming the radial and angular parameters into 3D
positions. The resulting Cartesian gripper configuration is then passed to the evaluation endpoint,
which executes the grasp evaluation procedure and returns the objective values.

This process can be written as
\begin{equation}
    \mathbf{p} = \Phi(\mathbf{q}),
    \qquad
    \mathbf{g}(\mathbf{q})
    =
    \texttt{API\_call}(\mathbf{p})
    =
    [D(\mathbf{p}),\;R_g(\mathbf{p}),\;A(\mathbf{p})].
\end{equation}
where $\Phi(\cdot)$ denotes the transformation from the reduced parameterization to Cartesian
suction cup coordinates.
\nomenclature{$\Phi(\mathbf{q})$}{Transformation from reduced decision variables to Cartesian suction cup coordinates}

In the implementation, this transformation is carried out programmatically in the evaluation
pipeline before grasp generation and objective computation. For a candidate radius and list of
angles, the suction cup coordinates are constructed as
\begin{equation}
    \mathbf{c}_i = (r\cos(\alpha_i),\; r\sin(\alpha_i),\; z_0),
    \qquad i=1,2,3.
\end{equation}

This architecture enables:
\begin{itemize}
    \item Separation of concerns between optimization logic and domain-specific evaluation
    \item Scalability through independent deployment of the evaluation service
    \item Reusability of the optimization framework for different problem domains
\end{itemize}

\subsubsection{Grasp Evaluation}
\label{sec:grasp_eval_sc}

The evaluation service receives a Cartesian gripper configuration $\mathbf{p} =
[\mathbf{c}_1,\mathbf{c}_2,\mathbf{c}_3]$ and computes the associated grasp quality measures through
the following procedure.

\begin{enumerate}
    \item \textbf{Candidate Grasp Generation and Feasibility Filtering}

    Suction grasp candidates are generated using a method developed by Fraunhofer IPA, following the automatic suction grasp generation procedure
    described by ~\cite{Kleeberger}. Points are uniformly sampled on the surface of the workpiece, and each sampled point defines a candidate grasp
    through its position and local surface normal. The local surface normal determines the approach direction of the suction cup.

    In the original method, each suction grasp candidate is tested for feasible seal formation using a quasi-static spring model of a suction cup, as
    introduced in ~\cite{SC_quasi_static}. If the grasp point is considered feasible, the candidate is further tested with a gripper collider model
    for collisions between the gripper and the object.

    In this thesis, the method was extended to a gripper with three suction cups. For each candidate gripper pose, the seal-formation test is
    performed for all three suction cup contact positions $\mathbf{c}_1$, $\mathbf{c}_2$, and $\mathbf{c}_3$. A candidate is considered feasible only
    if all three suction cups form a sufficient seal. Each feasible candidate is subsequently subjected to a geometric collision check: the gripper
    geometry is represented by a sampled point set, transformed into the workpiece frame by the candidate pose, and tested for proximity against the
    workpiece point cloud. A candidate is rejected if any transformed gripper point lies closer than the clearance threshold $\delta =
    0.5\,\mathrm{mm}$ to any workpiece point. The candidates that pass both the seal-formation and the collision checks form the set of valid grasps
    $\mathcal{G}_{\mathrm{valid}}$, which is the input to the subsequent clustering step.
    \nomenclature{$\mathbf{p}$}{Cartesian gripper configuration defined by the three suction cup
    contact positions (see also $\mathbf{p}$ for Cartesian suction cup configuration)}
    \nomenclature{$\mathbf{c}_i$}{Cartesian contact position of suction cup $i$, with $i \in \{1,2,3\}$ (see also $\mathbf{c}_i$ for suction cup positions)}

\item \textbf{Clustering of Grasps}

The set $\mathcal{G}_{\mathrm{valid}}$ often contains many grasps that are nearly
identical in position and orientation. When the number of valid grasps
$m = |\mathcal{G}_{\mathrm{valid}}|$ exceeds a threshold $m_{\max}$
(\texttt{max\_cluster}), a farthest-point sampling (FPS) procedure is applied to select
a diverse reduced subset. Otherwise, all valid grasps are used directly.

Each grasp $g_i$ is represented by a pose descriptor
\begin{equation}
    \mathbf{z}_i =
    [\,\mathrm{vec}(\mathbf{R}_i)^\top,\ \mathbf{t}_i^\top\,]^\top
    \in \mathbb{R}^{12},
\end{equation}
where $\mathbf{R}_i \in SO(3)$ is the orientation and $\mathbf{t}_i \in \mathbb{R}^3$
is the position. Translational and rotational differences are measured separately. The
translational distance
\begin{equation}
    d_{\mathrm{pos}}(i,j) = \|\mathbf{t}_i - \mathbf{t}_j\|_2
\end{equation}
is normalised by the diagonal extent of the axis-aligned bounding box of all candidate
positions,
\begin{equation}
    s_{\mathrm{pos}} =
    \left\lVert
        \max_i \mathbf{t}_i - \min_i \mathbf{t}_i
    \right\rVert_2 ,
    \qquad
    \hat{d}_{\mathrm{pos}}(i,j) = \frac{d_{\mathrm{pos}}(i,j)}{s_{\mathrm{pos}}},
\end{equation}
where the minimum and maximum are taken component-wise. In contrast to a linearised
comparison of rotation matrix entries, the orientation difference is evaluated as the
geodesic angle on $SO(3)$,
\begin{equation}
    d_{\mathrm{ori}}(i,j) =
    \arccos\!\left(
        \frac{\mathrm{tr}\!\left(\mathbf{R}_i \mathbf{R}_j^\top\right) - 1}{2}
    \right),
    \qquad
    \hat{d}_{\mathrm{ori}}(i,j) = \frac{d_{\mathrm{ori}}(i,j)}{\pi},
\end{equation}
with the argument of the arccosine numerically clamped to $[-1,1]$. Both normalised
terms are combined into a single pose metric using an orientation weight
$w \in [0,1]$ (\texttt{orientation\_weight}),
\begin{equation}
    \label{eq:combined-metric}
    d_w(i,j) =
    (1 - w)\,\hat{d}_{\mathrm{pos}}(i,j) + w\,\hat{d}_{\mathrm{ori}}(i,j),
\end{equation}
so that $w = 0$ reduces to pure position-based sampling and $w = 1$ to pure orientation
-based sampling. All distances are evaluated on the GPU, which makes the pairwise
evaluations tractable for large candidate sets.

Cluster representatives are selected greedily by farthest-point sampling in the combined
metric \eqref{eq:combined-metric}. The first representative is the grasp closest to the
centroid of all candidate positions,
\begin{equation}
    i_1 = \arg\min_i \left\lVert \mathbf{t}_i - \bar{\mathbf{t}} \right\rVert_2 ,
    \qquad
    \bar{\mathbf{t}} = \frac{1}{m} \sum_{i=1}^{m} \mathbf{t}_i ,
\end{equation}
rather than the most extreme grasp, because an edge-first initialisation tends to omit
centrally located grasps when the output budget is small. Each subsequent representative
maximises its minimum distance to the already selected ones,
\begin{equation}
    i_{k+1} = \arg\max_{i}\ \min_{1 \le j \le k} d_w(i, i_j).
\end{equation}

Once the number of globally sampled representatives reaches the requested output budget,
the remaining slots of the sampling pool are filled with orientation-diverse grasps from
local neighbourhoods. For a representative $i_k$, all unselected grasps within a spatial
radius $\delta_{\mathrm{intra}}$ (\texttt{intra\_cluster\_dist}) form the neighbourhood
\begin{equation}
    \mathcal{N}(i_k) =
    \left\{ i \;:\; d_{\mathrm{pos}}(i, i_k) < \delta_{\mathrm{intra}} \right\},
\end{equation}
from which up to $q_{\mathrm{intra}} - 1$ additional grasps
(\texttt{intra\_cluster\_num}) are chosen by a max--min selection on the geodesic
distance $d_{\mathrm{ori}}$. This captures distinct orientations at nearly identical
positions.

The procedure yields a candidate pool of $N_{\mathrm{pool}}$ representatives
(\texttt{fps\_pool\_size}), which is subsequently trimmed to the final subset size by an
elbow-based criterion. During sampling, the covering radius after $N$ selections,
\begin{equation}
    c_N = \max_{i}\ \min_{1 \le j \le N} d_w(i, i_j),
\end{equation}
is recorded as a coverage curve; $c_N$ is non-increasing in $N$ and quantifies how well
the selected subset represents the full candidate set. The final number of grasps is the
smallest $N$ within the admissible bounds $[N_{\min}, N_{\max}]$
(\texttt{auto\_n\_min}, \texttt{auto\_n\_max}) whose relative coverage gain drops below a
threshold $\tau$ (\texttt{auto\_n\_relative\_gain\_threshold}),
\begin{equation}
    N^{*} =
    \min \left\{
        N \in [N_{\min},\, N_{\max}-1] \;:\;
        \frac{c_N - c_{N+1}}{c_N} < \tau
    \right\},
\end{equation}
and $N^{*} = N_{\max}$ if no such $N$ exists. The bounds are clamped to the requested
maximum $m_{\max}$ and to the pool size, so that the selected count never exceeds the
available representatives. The resulting clustered subset
\begin{equation}
    \mathcal{G}_{\mathrm{clustered}}
    = \left\{ g_{i_1}, \dots, g_{i_{N^{*}}} \right\}
    \subseteq \mathcal{G}_{\mathrm{valid}}
\end{equation}
is then used for diversity evaluation.

\nomenclature{$\mathcal{G}_{\mathrm{clustered}}$}{Reduced clustered grasp set used for diversity evaluation}
\nomenclature{$\mathbf{z}_i$}{Pose descriptor of grasp $i$, composed of rotation and translation entries}
\nomenclature{$d_{\mathrm{pos}}(i,j)$}{Pairwise translational distance between grasp poses $i$ and $j$}
\nomenclature{$d_{\mathrm{ori}}(i,j)$}{Pairwise geodesic orientation distance between grasp poses $i$ and $j$}
\nomenclature{$s_{\mathrm{pos}}$}{Normalisation scale for translational distances (bounding-box diagonal)}
\nomenclature{$d_w(i,j)$}{Combined pose distance with orientation weight $w$}
\nomenclature{$w$}{Orientation weight in the combined pose distance}
\nomenclature{$m_{\max}$}{Maximum number of grasps retained by clustering}
\nomenclature{$\delta_{\mathrm{intra}}$}{Spatial radius for selecting additional local representatives}
\nomenclature{$q_{\mathrm{intra}}$}{Maximum number of grasps selected from one local neighbourhood}
\nomenclature{$N_{\mathrm{pool}}$}{Size of the farthest-point sampling candidate pool}
\nomenclature{$c_N$}{Covering radius of the selected subset after $N$ samples}
\nomenclature{$N_{\min}$, $N_{\max}$}{Lower and upper bound for the final number of clustered grasps}
\nomenclature{$\tau$}{Relative coverage-gain threshold of the elbow criterion}
\nomenclature{$N^{*}$}{Final number of clustered grasps selected by the elbow criterion}


\item \textbf{Diversity Metric Calculation}

A diversity score $D(\mathbf{p})$ is computed to quantify how broadly the resulting valid
grasp set is distributed over the object, both in position and in orientation.

Let the set used for diversity evaluation contain $n$ grasps. If clustering is performed,
then
\begin{equation}
    n = |\mathcal{G}_{\mathrm{clustered}}|,
\end{equation}
otherwise
\begin{equation}
    n = |\mathcal{G}_{\mathrm{valid}}|.
\end{equation}
For $n < 2$ no pairwise comparison is defined; in this case the corresponding component is
set to zero.

The evaluated positions are the projections of the grasp points onto the mesh surface,
\begin{equation}
    \mathbf{y}_i = \arg\min_{\boldsymbol{\xi} \in \partial \mathcal{M}}
    \|\boldsymbol{\xi} - \mathbf{x}_i\|_2 ,
    \qquad i = 1,\dots,n,
\end{equation}
where $\partial \mathcal{M}$ denotes the surface of the object mesh and $\mathbf{x}_i$ the
original grasp point. The associated orientation matrices are
$\{\mathbf{R}_1,\dots,\mathbf{R}_n\}$, with $\mathbf{R}_i \in SO(3)$.

To normalise the spatial terms with respect to object size, a characteristic object scale
is defined as the Euclidean norm of the axis-aligned bounding-box extents,
\begin{equation}
    s = \|\mathbf{e}\|_2 ,
\end{equation}
where $\mathbf{e} \in \mathbb{R}^3$ contains the bounding-box side lengths.

The spatial diversity is assessed by two complementary quantities. The first is the
normalised mean pairwise Euclidean distance between the grasp positions,
\begin{equation}
    D_{\mathrm{spread}}(\mathbf{p}) =
    \frac{1}{s} \cdot
    \frac{1}{n(n-1)}
    \sum_{\substack{i,j=1 \\ i \neq j}}^{n}
    \|\mathbf{y}_i - \mathbf{y}_j\|_2 .
\end{equation}
Because a large mean distance can also be achieved by grasps concentrated on two opposite
faces, a second term evaluates the angular coverage of the grasp positions as seen from
the object centroid $\mathbf{c}$. With the normalised view directions
\begin{equation}
    \hat{\mathbf{u}}_i =
    \frac{\mathbf{y}_i - \mathbf{c}}{\|\mathbf{y}_i - \mathbf{c}\|_2 + \varepsilon},
\end{equation}
where $\varepsilon$ is a small constant guarding against division by zero, the pairwise
cosine distance
\begin{equation}
    d_{\cos}(i,j) = 1 - \hat{\mathbf{u}}_i^\top \hat{\mathbf{u}}_j \in [0,2]
\end{equation}
vanishes for identical directions and attains its maximum for opposing directions. Its
normalised mean defines the coverage term
\begin{equation}
    D_{\mathrm{cov}}(\mathbf{p}) =
    \mathrm{clip}\!\left(
        \frac{1}{2} \cdot
        \frac{1}{n(n-1)}
        \sum_{\substack{i,j=1 \\ i \neq j}}^{n}
        d_{\cos}(i,j),
        \; 0, \; 1
    \right).
\end{equation}
Both contributions are combined using the angular-coverage weight
$\beta_{\mathrm{cov}} = 0.5$:
\begin{equation}
    D_{\mathrm{pos}}(\mathbf{p}) =
    (1-\beta_{\mathrm{cov}})\,D_{\mathrm{spread}}(\mathbf{p})
    +
    \beta_{\mathrm{cov}}\,D_{\mathrm{cov}}(\mathbf{p}).
\end{equation}
so that a high spatial score requires the grasps to be both mutually distant and
distributed around the object rather than clustered on a single side.

The orientation component is computed from the mean pairwise geodesic distance between
rotation matrices. For two grasp orientations
$\mathbf{R}_a, \mathbf{R}_b \in SO(3)$, the relative rotation is
\begin{equation}
    \Delta \mathbf{R} = \mathbf{R}_a^\top \mathbf{R}_b
\end{equation}
and the associated geodesic distance is
\begin{equation}
    d_{\mathrm{geo}}(\mathbf{R}_a,\mathbf{R}_b)
    = \cos^{-1}\!\left(
        \frac{\mathrm{tr}(\Delta \mathbf{R}) - 1}{2}
    \right),
\end{equation}
where the argument of $\cos^{-1}$ is clipped to $[-1,1]$ for numerical stability. The
normalised orientation diversity follows as
\begin{equation}
    D_{\mathrm{rot}}(\mathbf{p}) =
    \frac{1}{\pi} \cdot
    \frac{2}{n(n-1)}
    \sum_{1 \leq i < j \leq n}
    d_{\mathrm{geo}}(\mathbf{R}_i,\mathbf{R}_j).
\end{equation}

The final diversity score weights the spatial component more strongly than the orientation
component, since the orientation of suction grasps is largely determined by the local
surface normal and therefore varies less independently. With $\lambda = 0.65$,
\begin{equation}
    D(\mathbf{p}) =
    \mathrm{clip}\!\left(
        \lambda\, D_{\mathrm{pos}}(\mathbf{p}) +
        (1-\lambda)\, D_{\mathrm{rot}}(\mathbf{p}),
        \; 0, \; 1
    \right),
\end{equation}
where the final clipping prevents values marginally exceeding unity due to floating-point
rounding. For large grasp sets the orientation term saturates around
$D_{\mathrm{rot}} \approx 0.7$, as uniformly distributed rotations do not reach the
theoretical maximum of the normalised geodesic mean.

\nomenclature{$D(\mathbf{p})$}{Diversity metric of the grasp distribution (see also $D_{\mathrm{pos}}$, $D_{\mathrm{rot}}$ in theoretical background)}
\nomenclature{$D_{\mathrm{pos}}(\mathbf{p})$}{Spatial component of the diversity metric (see also $D_{\mathrm{pos}}$ in theoretical background)}
\nomenclature{$D_{\mathrm{spread}}(\mathbf{p})$}{Normalised mean pairwise distance between grasp positions}
\nomenclature{$D_{\mathrm{cov}}(\mathbf{p})$}{Angular coverage of grasp positions around the object centroid}
\nomenclature{$D_{\mathrm{rot}}(\mathbf{p})$}{Orientation component of the diversity metric (see also $D_{\mathrm{rot}}$ in theoretical background)}
\nomenclature{$d_{\mathrm{geo}}(\mathbf{R}_a,\mathbf{R}_b)$}{Geodesic distance between two grasp orientations (see also $d_{\mathrm{geo}}$ in theoretical background)}
\nomenclature{$d_{\cos}(i,j)$}{Cosine distance between the centroid-referenced directions of grasps $i$ and $j$}
\nomenclature{$\Delta\mathbf{R}$}{Relative rotation between two grasp orientations}
\nomenclature{$\hat{\mathbf{u}}_i$}{Normalised direction from the object centroid to grasp position $i$}
\nomenclature{$\mathbf{c}$}{Centroid of the object mesh}
\nomenclature{$\beta_{\mathrm{cov}}$}{Weight of the angular coverage term within the spatial diversity component}
\nomenclature{$\lambda$}{Weight of the spatial component in the total diversity score}
\nomenclature{$s$}{Characteristic object scale derived from the mesh bounding box (see also $s$ in theoretical background)}
\nomenclature{$\mathbf{e}$}{Axis-aligned bounding-box side lengths of the object (see also $\mathbf{e}$ in theoretical background)}
\nomenclature{$\mathbf{y}_i$}{Surface-projected position of grasp $i$ used in diversity evaluation (see also $\mathbf{y}_i$ in theoretical background)}
\nomenclature{$n$}{Number of grasps used for diversity evaluation (see also $n$ in theoretical background)}

\item \textbf{Normalised Grasp Count}

Besides the spatial distribution of the grasps, their absolute number is relevant as well,
since a layout that offers more valid grasps is more robust with respect to object pose and
surrounding scene constraints. The evaluated grasp count is taken from the same set that is
used for the diversity metric, i.e.
\begin{equation}
    n =
    \begin{cases}
        |\mathcal{G}_{\mathrm{clustered}}|, & \text{if clustering is performed},\\[2pt]
        |\mathcal{G}_{\mathrm{valid}}|,     & \text{otherwise}.
    \end{cases}
\end{equation}
To keep this quantity comparable with the remaining objectives, it is expressed relative to
the clustering budget $m_{\max}$ (\texttt{max\_cluster}), which is the maximum number of
grasps retained by the clustering stage,
\begin{equation}
    R_g(\mathbf{p}) = \frac{n}{m_{\max}} \in [0,1].
\end{equation}
Since clustering is only applied when more than $m_{\max}$ valid grasps exist, the condition
$n \leq m_{\max}$ holds in both cases and the ratio is bounded by unity without further
clipping. A value of one indicates that the layout produced at least as many valid grasps as
the clustering budget permits, whereas a value of zero corresponds to a layout for which no
valid grasp was found.
\nomenclature{$R_g(\mathbf{p})$}{Normalised grasp count, i.e.\ the number of valid grasps divided by the clustering budget $m_{\max}$}
\nomenclature{$m_{\max}$}{Maximum number of grasps retained by the clustering stage (\texttt{max\_cluster})}

\item \textbf{Triangle Area Calculation}

The area of the triangle formed by the three suction cups is computed from their Cartesian
coordinates. If the planar coordinates of the three cups are $(x_1,y_1)$, $(x_2,y_2)$ and
$(x_3,y_3)$, then
\begin{equation}
    A_{\mathrm{base}}(\mathbf{p}) =
    \frac{1}{2}
    \left|
        (x_2 - x_1)(y_3 - y_1) - (x_3 - x_1)(y_2 - y_1)
    \right|.
\end{equation}

To make this objective comparable in magnitude with the remaining objectives, the area is
normalised by the largest triangle area that can be realised within the admissible mounting
region. The suction-cup reference points are restricted to a circle of radius $r_{\max}$,
and the maximum-area triangle inscribed in a circle is the equilateral one, so that
\begin{equation}
    A_{\max} = \frac{3\sqrt{3}}{4}\, r_{\max}^{2}.
\end{equation}
The radius $r_{\max}$ is not a global constant but is determined per workpiece as the
largest circle that still fits inside the admissible mounting bounds. Consequently
$A(\mathbf{p})$ expresses the achieved area relative to what is geometrically possible for
that specific workpiece.

The normalised area objective is therefore
\begin{equation}
    A(\mathbf{p}) = \kappa \, \frac{A_{\mathrm{base}}(\mathbf{p})}{A_{\max}}
    \in [0, \kappa],
\end{equation}
where $\kappa$ is a scaling constant that is set to $\kappa = 1$ in this work, so that
$A(\mathbf{p}) \in [0,1]$.

This objective favours gripper layouts with a wider support polygon, which is associated
with improved grasp stability.
\nomenclature{$A_{\mathrm{base}}(\mathbf{p})$}{Unnormalised area of the triangle formed by the three suction cups}
\nomenclature{$A(\mathbf{p})$}{Normalised triangle area objective (see also triangle area for grasp stability)}
\nomenclature{$A_{\max}$}{Maximum triangle area attainable within the admissible mounting circle}
\nomenclature{$r_{\max}$}{Radius of the admissible mounting circle for the suction-cup reference points}
\nomenclature{$\kappa$}{Scaling constant of the normalised area objective}


\paragraph{Continuity considerations}
Clustering is only applied when the number of valid grasps exceeds the predefined threshold
$m_{\max}$ (\texttt{max\_cluster}). If fewer valid grasps are available, clustering is skipped and
the diversity metric is evaluated directly on the valid grasp set. Therefore, the omission of
clustering does not by itself introduce a discontinuity in the objective function.

A discontinuity may arise only in the degenerate case where fewer than two valid grasps remain after
collision filtering. Since both the spatial and orientation components of the diversity metric are
based on pairwise distances, these terms cannot be evaluated meaningfully for $|\mathcal{G}_{\mathrm{valid}}|<2$. In that case,
the corresponding terms are set to zero. As a result, small changes in the gripper configuration may
cause an abrupt transition between a nonzero diversity score and a degenerate zero-valued diversity
contribution.

To reduce the likelihood of this case, the number of initially sampled surface points was increased
from 100 to 1000. The denser surface sampling produces a larger set of candidate grasps and makes it
much less likely that fewer than two valid grasps remain after filtering.
\end{enumerate}
\subsubsection{Objective Evaluation}

After grasp evaluation, each candidate layout is assigned the natural objective vector
\begin{equation}
    \mathbf{g}(\mathbf{q})
    =
    [D(\mathbf{p}),\;R_g(\mathbf{p}),\;A(\mathbf{p})],
    \qquad
    \mathbf{p} = \Phi(\mathbf{q}).
\end{equation}
where the ordering corresponds to the return order of the evaluation pipeline.
\nomenclature{$\mathbf{g}(\mathbf{q})$}{Natural multi-objective function vector containing grasp diversity, normalised grasp count and normalised triangle area}


All three components are maximised:

\begin{itemize}
    \item $D(\mathbf{p})$ -- the diversity score, where larger values indicate that the
          valid grasps are more broadly distributed over the object in position and
          orientation.
    \item $R_g(\mathbf{p})$ -- the normalised grasp count, where larger values indicate
          that the layout can grasp the object in more distinct ways and therefore offers
          greater robustness with respect to object pose and scene constraints.
    \item $A(\mathbf{p})$ -- the normalised triangle area spanned by the three suction
          cups, where larger values correspond to a wider support polygon.
\end{itemize}

The objective-direction convention is handled differently by the two optimisation
frameworks. Since \texttt{pymoo} assumes minimisation, the evolutionary algorithms receive
the negated objective vector
\begin{equation}
    \mathbf{f}_{\mathrm{EA}}(\mathbf{q})
    =
    [-D(\mathbf{p}),\;-R_g(\mathbf{p}),\;-A(\mathbf{p})].
\end{equation}
\nomenclature{$\mathbf{f}_{\mathrm{EA}}(\mathbf{q})$}{Minimisation-form objective vector supplied to the evolutionary algorithms (see also $\mathbf{f}_{\mathrm{EA}}$ in two-finger optimization)}

In the Bayesian implementation, the values returned by the external evaluation service are
left unchanged, since \texttt{Ax} assumes maximisation,
\begin{equation}
    \mathbf{f}_{\mathrm{BO}}(\mathbf{q})
    =
    [D(\mathbf{p}),\;R_g(\mathbf{p}),\;A(\mathbf{p})].
\end{equation}
\nomenclature{$\mathbf{f}_{\mathrm{BO}}(\mathbf{q})$}{Objective observations supplied to Ax (see also $\mathbf{f}_{\mathrm{BO}}$ in two-finger optimization)}

The optimisation direction is specified through the Ax objective expressions, so the
observations themselves are not negated before being reported. This preserves the original
physical meaning of the objective values in the returned Pareto set.


\subsection{Optimization Loop for Multi-Suction Cup Gripper}

The candidate-generation process differs between evolutionary and Bayesian optimization, while the
external candidate-evaluation procedure remains identical.

\subsubsection{Evolutionary Optimization Loop}

The evolutionary optimization proceeds according to
\begin{equation}
    \mathcal{U}^{(g+1)}
    =
    \mathcal{A}
    \left(
        \mathcal{U}^{(g)},
        \mathbf{f}_{\mathrm{EA}}
    \right),
    \qquad
    g=0,1,\dots,
\end{equation}
where $\mathcal{A}$ denotes the selected evolutionary algorithm.
\nomenclature{$\mathcal{A}$}{Selected multi-objective evolutionary algorithm (see also $\mathcal{A}$ in two-finger optimization)}
\nomenclature{$g$}{Evolutionary generation index}

Depending on the selected algorithm, the new population is generated through operations such as
selection, crossover, mutation, decomposition, particle movement, or survival based on hypervolume
contribution. Every candidate in the resulting population is evaluated through the common API-based
grasp evaluation pipeline.

\subsubsection{Bayesian Optimization Loop}

The Bayesian optimization procedure maintains the observation dataset $\mathcal{D}_t$. During the
initialization phase, candidates are generated from the center of the search space and by Sobol
sampling. During the model-based phase, the available observations are used to fit probabilistic
surrogate models:
\begin{equation}
    \widehat{\mathbf{g}}_t
    =
    \mathcal{M}(\mathcal{D}_t),
\end{equation}
where $\mathcal{M}$ denotes the surrogate-modeling procedure used by Ax and BoTorch.
\nomenclature{$\mathcal{M}$}{Bayesian surrogate-modeling procedure (see also $\mathcal{M}$ in two-finger optimization)}
\nomenclature{$\widehat{\mathbf{g}}_t$}{Probabilistic surrogate approximation after $t$ evaluations (see also $\widehat{\mathbf{g}}_t$ in two-finger optimization)}

The surrogate models provide predictions of the objective values together with their uncertainty. An
acquisition function $\alpha(\mathbf{q}\mid\mathcal{D}_t)$ uses this information to balance the
exploration of uncertain regions and the exploitation of regions expected to improve the current
Pareto frontier.

For a requested batch size $b$, the next candidate batch is generated as
\begin{equation}
    \mathcal{Q}_{t+1}
    =
    \operatorname*{arg\,max}_{
        \substack{
            \mathcal{Q}\subset\mathcal{X}\\
            |\mathcal{Q}|\leq b
        }
    }
    \alpha(\mathcal{Q}\mid\mathcal{D}_t),
\end{equation}
where $\mathcal{X}$ is the bounded decision space. \nomenclature{$\mathcal{Q}_{t+1}$}{Batch of
candidates generated by Bayesian optimization (see also $\mathcal{X}_{t+1}$ in two-finger optimization for similar concept)}
\nomenclature{$\mathcal{X}$}{Bounded decision space (see also $\mathcal{X}$ in two-finger optimization)}
\nomenclature{$\alpha$}{Bayesian acquisition function (see also $\alpha$ in two-finger optimization)}

Each candidate $\mathbf{q}\in\mathcal{Q}_{t+1}$ is evaluated through the external evaluation
service. The resulting observation is then supplied to Ax as a completed trial:
\begin{equation}
    \mathcal{D}_{t+1}
    =
    \mathcal{D}_t
    \cup
    \left\{
        \left(
            \mathbf{q}_{t+1},
            \mathbf{g}(\mathbf{q}_{t+1})
        \right)
    \right\}.
\end{equation}

The surrogate model is subsequently updated using the expanded dataset, and the process is repeated
until the specified evaluation budget is exhausted. Although multiple candidates can be requested in
one batch, the current implementation submits the candidates in that batch sequentially to the
evaluation service.

\subsection{Termination Criteria for Multi-Suction Cup Gripper Optimization}

The termination criteria depend on the selected optimization approach.

\begin{table}[htbp]
\centering
\caption{Termination criteria for evolutionary and Bayesian optimization.}
\label{tab:termination_sc}
\begin{tabularx}{\textwidth}{@{}l l l >{\raggedright\arraybackslash}X@{}}
\toprule
\textbf{Approach} & \textbf{Criterion} & \textbf{Symbol} & \textbf{Description} \\
\midrule
Evolutionary & Design-space tolerance & $x_{\text{tol}}$ & Change in decision variables below
threshold \\

Evolutionary & Objective-space tolerance & $f_{\text{tol}}$ & Relative objective improvement below
threshold \\

Evolutionary & Maximum generations & $G_{\max}$ & Maximum number of generations reached \\

Evolutionary & Maximum evaluations & $E_{\max}$ & Maximum number of function evaluations reached \\

Evolutionary & Hypervolume & $HV_{\max}$ & Hypervolume progression stagnates  \\

Bayesian & Maximum batches & $B_{\max}$ & Maximum number of candidate-generation batches reached \\

Bayesian & Maximum evaluations & $E_{\max}^{\mathrm{BO}}$ & Bayesian evaluation budget exhausted \\
\bottomrule
\end{tabularx}
\end{table}

For evolutionary optimization, the design-space and objective-space tolerances are evaluated over a
sliding window of $w$ generations. \nomenclature{$x_{\text{tol}}$}{Design-space convergence
tolerance (see also $x_{\text{tol}}$ in two-finger optimization)} \nomenclature{$f_{\text{tol}}$}{Objective-space convergence tolerance
(see also $f_{\text{tol}}$ in two-finger optimization)}
\nomenclature{$G_{\max}$}{Maximum number of evolutionary generations (see also $G_{\max}$ in two-finger optimization)}
\nomenclature{$E_{\max}$}{Maximum number of evolutionary objective evaluations (see also $E_{\max}$ in two-finger optimization)}

In the current Bayesian implementation, termination is budget-based. The optimization performs at
most $B_{\max}$ candidate-generation batches, with up to $b$ candidates requested per batch.
Therefore,
\begin{equation}
    E_{\max}^{\mathrm{BO}} = B_{\max}b.
\end{equation}
No design-space or objective-space convergence tolerance is currently applied to the Bayesian
optimization loop.

\subsection{Performance Monitoring for Multi-Suction Cup Gripper}

For evolutionary optimization, the hypervolume indicator is tracked at each generation
via a callback mechanism:
\begin{equation}
    \mathrm{HV}^{(g)}
    =
    \mathrm{HV}\!\left(\mathcal{F}^{(g)},\mathbf{r}\right).
\end{equation}
\nomenclature{$\mathrm{HV}^{(g)}$}{Hypervolume indicator at evolutionary generation $g$}
\nomenclature{$\mathbf{r}$}{Reference point for hypervolume computation (see also $\mathbf{r}$ for reference point in theoretical background)}
\nomenclature{$\mathcal{F}^{(g)}$}{Objective vectors of the evolutionary population at generation $g$}

Here, $\mathbf{r} \in \mathbb{R}^{n_{\text{obj}}}$ is the reference point and
$\mathcal{F}^{(g)}$ denotes the set of objective vectors associated with the population
$\mathcal{U}^{(g)}$. An increasing hypervolume generally indicates progress toward a better approximation of
the Pareto front. Monotonicity is guaranteed only if the indicator is computed from a
cumulative nondominated archive; the hypervolume of the current population may fluctuate
between generations.

\subsection{Optional Hyperparameter Optimization for Multi-Suction Cup Gripper Problem}

The Evolutionary Optimization framework optionally supports automated tuning of algorithm
hyperparameters, such as crossover probability and mutation rate, using a meta-optimization
approach. A \texttt{HyperparameterProblem} is constructed around the selected algorithm and
evaluated over multiple random seeds. The mean hypervolume across runs serves as the meta-objective:
\begin{equation}
    \min_{\boldsymbol{\theta}}
    \left\{
        -\frac{1}{|\mathcal{Z}|}
        \sum_{s \in \mathcal{Z}} \text{HV}(\mathcal{F}_s,\mathbf{r})
    \right\}
\end{equation}
\nomenclature{$\boldsymbol{\theta}$}{Algorithm hyperparameters (see also $\boldsymbol{\theta}$ for hyperparameter optimization)}
\nomenclature{$\mathcal{Z}$}{Set of random seeds for meta-optimization (see also $\mathcal{Z}$ for random seeds)}

The best hyperparameters are then applied to the main optimization run.

\subsection{Output of Multi-Suction Cup Gripper Optimization}

The optimization produces a set of Pareto-optimal solutions:
\begin{equation}
    \mathcal{P}^*
    =
    \left\{
        \mathbf{q} \in \mathcal{X}
        \;\middle|\;
        \nexists\; \mathbf{q}' \in \mathcal{X}
        \text{ such that }
        \mathbf{g}(\mathbf{q}') \succeq \mathbf{g}(\mathbf{q})
        \text{ and }
        \mathbf{g}(\mathbf{q}') \neq \mathbf{g}(\mathbf{q})
    \right\}.
\end{equation}
\nomenclature{$\mathcal{P}^*$}{Set of Pareto-optimal solutions in the reduced decision space (see also $\mathcal{P}^{*}$ for Pareto-optimal solutions)}

The final output includes:
\begin{itemize}
    \item The Pareto-optimal set $\mathcal{P}^*$ in the reduced decision space, with corresponding
    objective values
    \item The associated radial gripper parameterizations $\mathbf{q} =
    [r,\alpha_1,\alpha_2,\alpha_3]$
    \item The hypervolume progression $\{\mathrm{HV}^{(0)},\mathrm{HV}^{(1)},\dots,\mathrm{HV}^{(G_{\mathrm{final}})}\}$
    \item Total optimization wall-clock time
\end{itemize}

\nomenclature{$G_{\mathrm{final}}$}{Final evolutionary generation reached by the optimization}

These solutions represent optimal trade-offs among grasp diversity, the number of valid grasps,
and triangle area, corresponding to grasp versatility, robustness, and stability.